\documentclass[12pt]{article}
\usepackage[utf8]{inputenc}
\usepackage[T1]{fontenc}
\usepackage[margin=1in]{geometry}
\usepackage{setspace}
\usepackage{xcolor}
\usepackage[colorlinks=true, urlcolor=blue!60!black, linkcolor=blue!60!black]{hyperref}
\setlength{\parindent}{2em}
\setlength{\parskip}{1.25em}
\renewcommand{\baselinestretch}{1.0}
\begin{document}
\begin{center}
{\LARGE \textit{Trump v. Barbara}, 609 U.S. \_\_\_ (2026) (No. 25-365)\par}
\vspace{0.8em}
{\large Case Brief --- Four-Page Edition\par}
\end{center}
\vspace{1em}

\section{Memory Jogger}

In a 6--3 judgment with a 5--4 constitutional majority, the Court invalidated Executive Order No.~14160 and held that the Citizenship Clause guarantees birthright citizenship to children born in the United States to parents who are unlawfully or temporarily present (\href{https://www.law.cornell.edu/supremecourt/text/25-365}{Cornell LII}).

\section{Detailed Case Facts}

On January 20, 2025, President Trump issued Executive Order No.~14160, declaring that citizenship does not extend to U.S.-born children whose mother was unlawfully present, or lawfully but temporarily present (as on a tourist, student, or work visa), and whose father was not a citizen or lawful permanent resident. The Order asserted that such children are not ``subject to the jurisdiction'' of the United States under the Fourteenth Amendment or 8 U.S.C.\ \S\ 1401(a), which codifies the same text. After \textit{Trump v.\ CASA, Inc.} (2025) curtailed universal injunctions, the ACLU filed a putative class action the same day on behalf of parents, led by a Honduran national suing as ``Barbara,'' and their children.

\section{Procedural History}

Judge Joseph Laplante of the District of New Hampshire certified a nationwide Rule 23(b)(2) class of children born on or after February 20, 2025, and issued a class-wide preliminary injunction. The Government appealed and sought certiorari before judgment, which the Supreme Court granted in December 2025. After argument on April 1, 2026, the Court affirmed on June 30, 2026.

\section{Judicial Votes}

Affirmed, \textbf{6--3} in judgment; \textbf{5--4} on the constitutional question.
\begin{itemize}
\item \textbf{Majority:} Chief Justice Roberts, joined by Justices Sotomayor, Kagan, Barrett, and Jackson.
\item \textbf{Concurring:} Justice Jackson, joined in part by Justice Sotomayor.
\item \textbf{Concurring in the judgment and dissenting in part:} Justice Kavanaugh (statutory grounds).
\item \textbf{Dissenting:} Justice Thomas, joined by Justice Gorsuch; Justices Alito and Gorsuch separately.
\end{itemize}

\section{Holding}

Children born in the United States to unlawfully or temporarily present parents are ``subject to the jurisdiction'' of the United States and citizens at birth. The phrase carries its ordinary territorial-sovereignty meaning, incorporating common-law \textit{jus soli} with only its narrow exceptions (children of diplomats, of hostile occupying forces, and historically of quasi-sovereign tribes). Neither the President nor Congress may deny citizenship the Clause guarantees. The majority did not rest on the statute, 8 U.S.C.\ \S\ 1401(a), and it read \textit{Wong Kim Ark}'s references to domicile as factual description, not limitation.

\section{Analysis of Opinions}

\textbf{Majority (Roberts).} Under \textit{Calvin's Case} and Blackstone, a person born within the dominions of the sovereign became a subject, however ``momentary and uncertain'' the parents' stay. That rule prevailed in the States and was reaffirmed in \textit{Lynch v.\ Clarke}. After \textit{Dred Scott} substituted a racial bloodline rule, the Civil Rights Act of 1866 and the Fourteenth Amendment restored \textit{jus soli}. In 1868, ``jurisdiction'' meant the power to govern within a territory; under \textit{The Schooner Exchange}, foreign visitors are ``amenable to the jurisdiction of the country.'' Reaffirming \textit{Wong Kim Ark}, Roberts rejected domicile and ``primary allegiance'' theories as ``dramatically revisionist.''

\textbf{Jackson (concurring).} Senator Trumbull withdrew a draft limited to ``persons of African descent'' for universal text, and Congress knowingly covered children of Chinese and Romani immigrants. Thomas's group-specific reading sits uneasily with a ``colorblind'' Constitution.

\textbf{Kavanaugh (concurring in the judgment).} Section 1401(a), enacted after \textit{Wong Kim Ark}, incorporated its rule, confirmed by 85 years of practice, so the President could not add exceptions. But the exceptions are not ``trapped in amber,'' and Congress could exclude such children.

\textbf{Thomas, with Gorsuch (dissenting).} The Clause constitutionalized the 1866 Act's requirement of birth plus \textit{domicile}. \textit{Calvin's Case} is a ``feudal'' relic, and \textit{Wong Kim Ark} covers only domiciled parents, so the facial challenge fails.

\textbf{Alito and Gorsuch (dissenting).} Alito called the decision ``a serious mistake'' that rewards ``birth tourists''; Gorsuch stressed that the Order was valid at least as to temporary visitors.

\section{Examples: Future Applications}

\begin{enumerate}
\item \textbf{Same-Side: Congressional Exclusion.} Congress amends \S\ 1401(a) to deny citizenship to children of parents who entered without inspection. The Clause itself guarantees citizenship, so only an amendment could change it.
\item \textbf{Same-Side: State Birth Certificates.} A State stamps ``Non-Citizen'' on certificates of children of tourist or student visa holders. Temporary visitors are amenable to U.S. jurisdiction, so the children are citizens.
\item \textbf{Opposite-Side: Ambassadors' Child.} A child born in Washington to accredited ambassadors is not a citizen, because diplomats are exempt from territorial jurisdiction.
\item \textbf{Opposite-Side: Hostile Occupation.} A child born to occupying forces on a seized American island is not born ``subject to the jurisdiction'' of the United States.
\item \textbf{Fence-Sitter: American Samoa.} A child born in an unincorporated territory is subject to U.S. jurisdiction, but whether the territory is ``in the United States'' under the \textit{Insular Cases}, which \textit{Barbara} did not overrule, is unresolved.
\end{enumerate}

\section{Critique}

\textbf{Originalist views.} Randy Barnett, Michael Ramsey, Evan Bernick, and the Cato Institute supported the majority's reading of original public meaning. Critics drawing on Kurt Lash and Ilan Wurman argued that the majority read ``jurisdiction'' out of the Clause and substituted feudal subjecthood for republican consent; post-liberal critics said the 5--4 split exposed originalism's indeterminacy.

\textbf{Progressive views.} Amanda Frost and Rachel Rosenbloom noted that four Justices signaled openness to congressional restrictions, so the guarantee held by one vote, and urged defenders to build on Jackson's antisubordination concurrence. Scholars also highlighted Rule 23(b)(2) nationwide classes as a remedial backstop after \textit{CASA}.

\section{Key Quotations}

\begin{quote}
``The Citizenship Clause uses jurisdiction in its ordinary sense---referring to the power of the United States to govern those within its territory.''
\end{quote}
\noindent --- Chief Justice Roberts, Opinion of the Court.

\begin{quote}
``Words appearing frequently in the Executive Order---`mother,' `father,' `lawful,' `temporary'---are absent from the Clause. For a simple reason: they did not matter.''
\end{quote}
\noindent --- Chief Justice Roberts, Opinion of the Court.

\begin{quote}
``The Reconstruction Amendments were an anticaste, antisubordination reset for the Nation, not a mere spot treatment for the dark stain of slavery.''
\end{quote}
\noindent --- Justice Jackson, concurring.

\end{document}
